\chapter{Ngô Bào Châu --- Fields Medal 2010}\label{ch2010ngobaochau:nguxf4-buxe0o-chuxe2u-fields-medal-2010-a-comprehensive-report-on-his-work-and-impact}

\begin{center}\rule{0.5\linewidth}{0.5pt}\end{center}

\textbf{Prepared:} July 2026 \textbf{Subject:} Ngô Bào Châu (b. June 28,
1972, Hanoi, Vietnam) \textbf{Position:} Professor, Department of
Mathematics, University of Chicago; Scientific Director, Vietnam
Institute for Advanced Study in Mathematics \textbf{Doctoral Advisor:}
Gérard Laumon \textbf{Fields Medal Citation:} \emph{``For his proof of
the Fundamental Lemma in the theory of automorphic forms through the
introduction of new algebro-geometric methods.''}

\begin{center}\rule{0.5\linewidth}{0.5pt}\end{center}

\section{Table of Contents}\label{ch2010ngobaochau:table-of-contents}

\begin{enumerate}
\def\labelenumi{\arabic{enumi}.}
\tightlist
\item
 Introduction
\item
 Biographical Background and Education
\item
 The Langlands Program: Historical Context
\item
 What Is the Fundamental Lemma?
\item
 The Fundamental Lemma Before Ngô
\item
 Geometric Methods: Hitchin Fibrations
\item
 Affine Springer Fibers
\item
 Ngô's Proof Strategy
\item
 The Main Theorem and Its Consequences
\item
 Applications to Number Theory
\item
 The Goormaghtigh Conjecture
\item
 Subsequent Developments in the Langlands Program
\item
 The Vietnam Institute for Advanced Study in Mathematics
\item
 Ngô's Broader Mathematical Influence
\item
 Awards and Recognition
\item
 Selected Major Publications
\item
 Ongoing Work and Future Directions
\item
 Conclusion
\item
 References and Further Reading
\end{enumerate}

\begin{prereqs} Algebraic geometry (stacks, intersection theory), number theory (automorphic forms, L-functions), representation theory. \end{prereqs}

\begin{objectives} Explain the fundamental lemma and its proof; describe the trace formula and its role in the Langlands program; understand the geometric interpretation of orbital integrals. \end{objectives}


\begin{center}\rule{0.5\linewidth}{0.5pt}\end{center}

\section{1. Introduction}\label{ch2010ngobaochau:introduction}

Ngô Bào Châu was awarded the Fields Medal in 2010 for his proof of the
Fundamental Lemma of the Langlands program --- one of the important
and long-standing problems in modern mathematics. The Fundamental Lemma,
formulated by Robert Langlands in the late 1970s and refined by
Shelstad, Kottwitz, and others, was a technical conjecture that stood as
a critical obstacle to progress in the Langlands program. Ngô's proof,
which combined deep insights from algebraic geometry, representation
theory, and harmonic analysis, was immediately recognized as a
monumental achievement.

Ngô is the first Vietnamese mathematician to win the Fields Medal. His
path to this achievement was notable: born in Hanoi during the
Vietnam War, he was educated in Vietnam, Hungary, and France, and his
work reflects a unique blend of mathematical cultures.

\begin{center}\rule{0.5\linewidth}{0.5pt}\end{center}

\section{2. Biographical Background and
Education}\label{ch2010ngobaochau:biographical-background-and-education}

\subsection{2.1 Early Life in Vietnam}\label{ch2010ngobaochau:early-life-in-vietnam}

Ngô Bào Châu was born on June 28, 1972, in Hanoi, North Vietnam, just as
the Vietnam War was entering its final years. His father, Ngô Huy Cẩn,
was a physicist and professor at the Vietnam National University, Hanoi,
specializing in solid-state physics. His mother, Trần Lưu Vân Hiên, was
a medical doctor.

Growing up in post-war Vietnam, Ngô showed exceptional mathematical
talent from an early age. He was accepted into the Mathematics and
Informatics Magnet High School of Hanoi University of Science, a
specialized school for mathematically gifted students in Vietnam.

\subsection{2.2 Mathematical Olympiads}\label{ch2010ngobaochau:mathematical-olympiads}

Ngô's mathematical talent was internationally recognized through his
participation in the International Mathematical Olympiad (IMO). He
competed twice:

\begin{itemize}
\tightlist
\item
 \textbf{1988 (Canberra, Australia):} Gold medal (the only Vietnamese
 contestant to receive a perfect score of 42/42)
\item
 \textbf{1989 (Braunschweig, West Germany):} Gold medal
\end{itemize}

These achievements earned him a scholarship to study in Hungary.

\subsection{2.3 University Education}\label{ch2010ngobaochau:university-education}

Ngô studied at the Eötvös Loránd University in Budapest, Hungary, from
1989 to 1992. However, due to financial difficulties and the changing
political landscape in Eastern Europe, he transferred to the Université
Paris-Sud in Orsay, France, in 1992.

He completed his undergraduate studies at Paris VI (Pierre and Marie
Curie University) and his master's degree at Université Paris-Sud.

\subsection{2.4 Doctoral Studies}\label{ch2010ngobaochau:doctoral-studies}

Ngô earned his PhD in 1997 at Université Paris-Sud under the supervision
of Gérard Laumon, a leading expert in the Langlands program. His thesis,
``Le lemme fondamental de Jacquet et Ye,'' already addressed aspects of
the Fundamental Lemma, focusing on the Jacquet--Ye variant for unitary
groups.

This work established Ngô as a rising star in the Langlands program. His
approach --- using algebraic geometry and especially the geometry of
affine Springer fibers --- would later become central to his complete
proof of the Fundamental Lemma.

\subsection{2.5 Academic Career}\label{ch2010ngobaochau:academic-career}

After his PhD, Ngô held positions at: - \textbf{Université Paris-Sud
(Orsay):} Assistant Professor (1998--2003) - \textbf{Institut des Hautes
Études Scientifiques (IHÉS):} Professor (2003--2005) -
\textbf{Université Paris-Sud (Orsay):} Professor (2005--2007) -
\textbf{Institute for Advanced Study, Princeton:} Professor (2007--2011)
- \textbf{University of Chicago:} Professor (2011--present)

He also serves as the Scientific Director of the Vietnam Institute for
Advanced Study in Mathematics (VIASM), where he has played a crucial
role in developing mathematics in Vietnam.

\begin{center}\rule{0.5\linewidth}{0.5pt}\end{center}

\section{3. The Langlands Program: Historical
Context}\label{ch2010ngobaochau:the-langlands-program-historical-context}

\subsection{3.1 The Vision}\label{ch2010ngobaochau:the-vision}

The Langlands program, initiated by Robert Langlands in 1967, is a vast
web of conjectures connecting number theory, representation theory, and
harmonic analysis. At its core is the idea that the absolute Galois
group of a number field and the automorphic representations of certain
algebraic groups are closely related.

The program has been described as ``a grand unified theory of
mathematics'' --- it connects seemingly unrelated areas such as: -
Diophantine equations (number theory) - Modular forms (complex analysis)
- Representations of Lie groups (algebra) - Spectral theory (analysis)

\subsection{3.2 Key Components}\label{ch2010ngobaochau:key-components}

The Langlands program includes several major components: -
\textbf{Reciprocity:} Connecting Galois representations to automorphic
forms (generalizing class field theory) - \textbf{Functoriality:}
Transfer of automorphic representations between different groups -
\textbf{Endoscopy:} A framework for comparing automorphic
representations of different groups, where the Fundamental Lemma plays a
crucial role

\subsection{3.3 The Role of Endoscopy}\label{ch2010ngobaochau:the-role-of-endoscopy}

Endoscopy, developed by Langlands, Shelstad, Kottwitz, and Arthur, is a
powerful technique for understanding automorphic representations of
reductive groups. It relates representations of a given group to those
of its ``endoscopic groups'' --- smaller groups that are closely
related.

The Fundamental Lemma asserts a key identity between orbital integrals
on a group and those on its endoscopic groups. Without this lemma, the
endoscopic transfer could not be established, and the Langlands program
could not progress.

\begin{center}\rule{0.5\linewidth}{0.5pt}\end{center}

\section{4. What Is the Fundamental
Lemma?}\label{ch2010ngobaochau:what-is-the-fundamental-lemma}

\subsection{4.1 Orbital Integrals}\label{ch2010ngobaochau:orbital-integrals}

Orbital integrals are integrals over conjugacy classes in a Lie group.
For a reductive group \(G\) over a local field \(F\), and a function
\(f\) on \(G(F)\), the orbital integral at a semisimple element
\(\gamma \in G(F)\) is:

\[
O_\gamma(f) = \int_{G_\gamma(F) \backslash G(F)} f(g^{-1}\gamma g) \, dg
\]

where \(G_\gamma\) is the centralizer of \(\gamma\). These integrals
appear naturally in the trace formula, an effective tool for studying
automorphic representations.

\subsection{4.2 The Statement}\label{ch2010ngobaochau:the-statement}

The Fundamental Lemma relates orbital integrals on a reductive group
\(G\) to orbital integrals on its endoscopic groups \(H\). Roughly
speaking, it states that for matching functions \(f\) on \(G\) and
\(f^H\) on \(H\) (related by a transfer of orbital integrals), the
orbital integrals on \(G\) and \(H\) satisfy:

\[
O_\gamma(f) = \sum_{H} \iota(G, H) \cdot O_{\gamma_H}(f^H)
\]

where \(\iota(G, H)\) are constants.

More precisely, for a local field \(F\) of characteristic zero (or large
positive characteristic), and for elements \(\gamma \in G(F)\) and
\(\gamma_H \in H(F)\) that are ``stably conjugate,'' the orbital
integrals of the unit function of the hyperspecial maximal compact
subgroup satisfy:

\[
O_\gamma(\mathbf{1}_{K}) = \sum_{\kappa} \iota(G, H_\kappa) \cdot O_{\gamma_\kappa}(\mathbf{1}_{K_\kappa})
\]

where the sum runs over endoscopic groups.

\subsection{4.3 Why It Is Fundamental}\label{ch2010ngobaochau:why-it-is-fundamental}

The lemma is called ``fundamental'' because it is required for: -
\textbf{The Arthur--Selberg trace formula:} To make the trace formula
work for comparing different groups - \textbf{Endoscopic transfer:} To
transfer automorphic representations between groups - \textbf{The
Langlands functoriality conjecture:} One of the central conjectures of
the Langlands program

Without the Fundamental Lemma, endoscopic transfer could not be
established, and large parts of the Langlands program would remain
conjectural.

\begin{center}\rule{0.5\linewidth}{0.5pt}\end{center}

\section{5. The Fundamental Lemma Before
Ngô}\label{ch2010ngobaochau:the-fundamental-lemma-before-nguxf4}

\subsection{5.1 Early History}\label{ch2010ngobaochau:early-history}

The Fundamental Lemma was first formulated by Langlands in the late
1970s. He quickly realized that many of his conjectures --- including
the functoriality conjecture --- depended on this lemma.

\subsection{5.2 Partial Results}\label{ch2010ngobaochau:partial-results}

Before Ngô's breakthrough, the Fundamental Lemma had been proved in a
number of special cases: - \textbf{Shelstad (1982):} Proved the lemma
for real and complex groups (archimedean case) - \textbf{Laumon (1987):}
Proved it for unitary groups in three variables - \textbf{Hales (1995):}
Proved it for \(\mathrm{SL}_n\) and \(\mathrm{Sp}_4\) -
\textbf{Waldspurger (1995--2006):} Made major advances, reducing the
problem to the case of positive characteristic - \textbf{Laumon and Ngô
(2004):} Proved the Fundamental Lemma for unitary groups

\subsection{5.3 The Reduction to Positive
Characteristic}\label{ch2010ngobaochau:the-reduction-to-positive-characteristic}

A key insight, due to Waldspurger, was that the Fundamental Lemma for
local fields of characteristic zero could be reduced to the case of
local fields of positive characteristic (function fields). This
reduction used sophisticated arguments from model theory and algebraic
geometry.

The problem thus became: prove the Fundamental Lemma for function
fields, i.e., for fields of the form \(\mathbb{F}_q((t))\).

\subsection{5.4 The Geometric
Formulation}\label{ch2010ngobaochau:the-geometric-formulation}

For function fields, the orbital integrals in the Fundamental Lemma can
be interpreted geometrically as counting points on certain algebraic
varieties --- \textbf{affine Springer fibers}. This geometric
interpretation, pioneered by Kazhdan and Lusztig (1988) and further
developed by Laumon, Ngô, and others, transformed the Fundamental Lemma
from an analytic statement into a geometric one.

\begin{center}\rule{0.5\linewidth}{0.5pt}\end{center}

\section{6. Geometric Methods: Hitchin
Fibrations}\label{ch2010ngobaochau:geometric-methods-hitchin-fibrations}

\subsection{6.1 The Hitchin Fibration}\label{ch2010ngobaochau:the-hitchin-fibration}

The Hitchin fibration, introduced by Nigel Hitchin in the context of
Higgs bundles on Riemann surfaces, is a map from the moduli space of
Higgs bundles to a vector space (the Hitchin base). For a reductive
group \(G\) over a curve \(X\), the Hitchin map takes a Higgs bundle
\((E, \theta)\) to the characteristic polynomial of \(\theta\).

Ngô realized that the Hitchin fibration provides the perfect geometric
framework for proving the Fundamental Lemma. The key is that the Hitchin
fibration is \textbf{proper} and its fibers are closely related to
affine Springer fibers.

\subsection{6.2 The Global-to-Local
Principle}\label{ch2010ngobaochau:the-global-to-local-principle}

Ngô's key insight was to work \textbf{globally} --- on a smooth
projective curve over a finite field --- rather than locally. The
Fundamental Lemma is a local statement, but Ngô proved it by studying
global objects (Higgs bundles) and then localizing.

This is analogous to the way that one proves local statements in number
theory by embedding them in global fields and using global methods.

\subsection{6.3 The Support Theorem}\label{ch2010ngobaochau:the-support-theorem}

A crucial technical tool in Ngô's proof is the \textbf{support theorem}
for the cohomology of the Hitchin fibration. This theorem describes the
support of the direct image sheaves in terms of the Hitchin base and
provides the key ingredient for the endoscopic classification of orbital
integrals.

\begin{center}\rule{0.5\linewidth}{0.5pt}\end{center}

\section{7. Affine Springer Fibers}\label{ch2010ngobaochau:affine-springer-fibers}

\subsection{7.1 Definition}\label{ch2010ngobaochau:definition}

For a reductive group \(G\) over a local field \(F\), and a regular
semisimple element \(\gamma \in G(F)\), the affine Springer fiber is the
space of \(G\)-valued points on the formal disk that fix \(\gamma\)
under conjugation. More precisely, letting \(K = F[[t]]\) be the ring of
integers of \(F\), the affine Springer fiber is:

\[
\mathcal{X}_\gamma = \{ g \in G(K) \backslash G(F) : g^{-1}\gamma g \in G(K) \}
\]

This is an infinite-dimensional space, but it can be approximated by
finite-dimensional varieties.

\subsection{7.2 Counting Points}\label{ch2010ngobaochau:counting-points}

The number of points on an affine Springer fiber over a finite field is
related to the orbital integral of the unit function:

\[
| \mathcal{X}_\gamma(\mathbb{F}_q) | = O_\gamma(\mathbf{1}_K)
\]

This connection --- between orbital integrals and point counts on
algebraic varieties --- is the key to Ngô's approach.

\subsection{7.3 Geometric Structure}\label{ch2010ngobaochau:geometric-structure}

Ngô, building on work of Kazhdan--Lusztig, Laumon, and others, developed
a deep understanding of the geometry of affine Springer fibers: -
\textbf{Purity:} The cohomology of affine Springer fibers is pure (in
the sense of Deligne's Weil conjectures). - \textbf{Dimensionality:} The
fibers have a well-defined dimension that can be computed in terms of
the root system. - \textbf{Stalks:} The stalks of the intersection
cohomology complex on the Hitchin base are given by the cohomology of
affine Springer fibers.

\begin{center}\rule{0.5\linewidth}{0.5pt}\end{center}

\section{8. Ngô's Proof Strategy}\label{ch2010ngobaochau:nguxf4s-proof-strategy}

\subsection{8.1 The Plan}\label{ch2010ngobaochau:the-plan}

Ngô's proof of the Fundamental Lemma proceeds in several major steps:

\subsection{8.2 Step 1: Reduction to the Geometric
Statement}\label{ch2010ngobaochau:step-1-reduction-to-the-geometric-statement}

Following Waldspurger, reduce the Fundamental Lemma for all local fields
to the case of function fields \(\mathbb{F}_q((t))\). Then reformulate
the lemma as a statement about point counts on affine Springer fibers.

\subsection{8.3 Step 2: The Hitchin
Fibration}\label{ch2010ngobaochau:step-2-the-hitchin-fibration}

Consider the moduli stack of \(G\)-Higgs bundles on a smooth projective
curve. The Hitchin fibration maps this stack to the Hitchin base. Study
its cohomology using the support theorem.

\subsection{8.4 Step 3: The Support
Theorem}\label{ch2010ngobaochau:step-3-the-support-theorem}

Prove that the direct image of the intersection cohomology complex under
the Hitchin fibration is a constant sheaf on the Hitchin base (after
appropriate base change). This is the \textbf{support theorem}, the
hardest technical part of the proof.

\subsection{8.5 Step 4: Endoscopic
Decomposition}\label{ch2010ngobaochau:step-4-endoscopic-decomposition}

Decompose the cohomology of the Hitchin fibration according to
endoscopic groups. The support theorem shows that the contributions of
different endoscopic groups are independent, which implies the
Fundamental Lemma.

\subsection{8.6 Step 5: Conclusion}\label{ch2010ngobaochau:step-5-conclusion}

The equality of orbital integrals --- the Fundamental Lemma --- follows
from the equality of point counts on the corresponding affine Springer
fibers, which follows from the endoscopic decomposition of the Hitchin
cohomology.

\begin{center}\rule{0.5\linewidth}{0.5pt}\end{center}

\section{9. The Main Theorem and Its
Consequences}\label{ch2010ngobaochau:the-main-theorem-and-its-consequences}

\subsection{9.1 Ngô's Theorem}\label{ch2010ngobaochau:nguxf4s-theorem}

In his 2008 paper ``Le lemme fondamental pour les algèbres de Lie'' (the
fundamental lemma for Lie algebras), Ngô proved the Fundamental Lemma
for \textbf{Lie algebras} (the so-called ``Lie algebra version''). The
group version was proved shortly after by establishing the equivalence
between the Lie algebra and group versions.

\subsection{9.2 Immediate Consequences}\label{ch2010ngobaochau:immediate-consequences}

Ngô's proof immediately implied: - \textbf{Endoscopic transfer:} The
transfer of orbital integrals between groups, which is the foundation of
endoscopic theory. - \textbf{The Arthur trace formula:} Completion of
the stabilization of the trace formula for classical groups. -
\textbf{The Langlands functoriality conjecture:} For many cases (e.g.,
from classical groups to general linear groups).

\subsection{9.3 Completing the Langlands
Program}\label{ch2010ngobaochau:completing-the-langlands-program}

Ngô's proof was described by Langlands himself as ``a great
achievement'' that ``removes the principal remaining obstacle to the
stabilization of the trace formula.'' The completion of the Fundamental
Lemma meant that: - The Arthur--Selberg trace formula could be fully
stabilized. - Endoscopic transfer could be established for all reductive
groups. - Major cases of the functoriality conjecture could be proved.

\begin{center}\rule{0.5\linewidth}{0.5pt}\end{center}

\section{10. Applications to Number
Theory}\label{ch2010ngobaochau:applications-to-number-theory}

\subsection{10.1 Automorphic Forms}\label{ch2010ngobaochau:automorphic-forms}

The most immediate application of Ngô's work is to the theory of
automorphic forms. The Fundamental Lemma is essential for proving that
automorphic representations transfer between groups, which is the core
of the Langlands program.

\subsection{10.2 Galois Representations}\label{ch2010ngobaochau:galois-representations}

The Langlands program connects automorphic forms to Galois
representations. The Fundamental Lemma, by enabling endoscopic transfer,
allows mathematicians to: - Construct Galois representations attached to
automorphic forms - Prove cases of the Ramanujan--Petersson conjecture -
Obtain new results on the distribution of eigenvalues of Hecke operators

\subsection{10.3 The Birch and Swinnerton-Dyer
Conjecture}\label{ch2010ngobaochau:the-birch-and-swinnerton-dyer-conjecture}

While Ngô's work does not directly prove the Birch and Swinnerton-Dyer
conjecture, it provides tools that may be relevant: - The Langlands
program is closely connected to the study of elliptic curves and their
L-functions. - Results on automorphic forms enabled by the Fundamental
Lemma can be used to study the arithmetic of elliptic curves.

\begin{center}\rule{0.5\linewidth}{0.5pt}\end{center}

\section{11. The Goormaghtigh
Conjecture}\label{ch2010ngobaochau:the-goormaghtigh-conjecture}

\subsection{11.1 The Conjecture}\label{ch2010ngobaochau:the-conjecture}

The Goormaghtigh conjecture (or Goormaghtigh's equation) states that the
only solutions to the exponential Diophantine equation:

\[
\frac{x^m - 1}{x - 1} = \frac{y^n - 1}{y - 1}
\]

with \(x > y > 1\) and \(m, n > 2\) are \((x, y, m, n) = (5, 2, 3, 3)\)
and \((x, y, m, n) = (90, 2, 3, 4)\).

This conjecture concerns repunits --- numbers whose decimal
representation consists entirely of 1's.

\subsection{11.2 Ngô's Contribution}\label{ch2010ngobaochau:nguxf4s-contribution}

In addition to his work on the Fundamental Lemma, Ngô has made
contributions to the Goormaghtigh conjecture, proving results about the
finiteness of solutions under certain conditions using methods from
Diophantine geometry.

\subsection{11.3 Significance}\label{ch2010ngobaochau:significance}

While less famous than the Fundamental Lemma, Ngô's work on the
Goormaghtigh conjecture demonstrates the breadth of his mathematical
interests --- from the most abstract Langlands program to concrete
Diophantine equations.

\begin{center}\rule{0.5\linewidth}{0.5pt}\end{center}

\section{12. Subsequent Developments in the Langlands
Program}\label{ch2010ngobaochau:subsequent-developments-in-the-langlands-program}

\subsection{12.1 The Geometric Langlands
Program}\label{ch2010ngobaochau:the-geometric-langlands-program}

The geometric Langlands program, which replaces number fields with
function fields of curves over complex numbers, has been one of the most
active areas of mathematics since the 1990s. Ngô's methods ---
especially his use of the Hitchin fibration --- have been directly
influential in geometric Langlands.

\subsection{12.2 The Work of Laurent Lafforgue and Vincent
Lafforgue}\label{ch2010ngobaochau:the-work-of-laurent-lafforgue-and-vincent-lafforgue}

Laurent Lafforgue (Fields Medal 2002) proved the Langlands
correspondence for \(\mathrm{GL}_n\) over function fields. Vincent
Lafforgue (Breakthrough Prize 2019) generalized this to all reductive
groups. Ngô's proof of the Fundamental Lemma was a crucial ingredient in
Vincent Lafforgue's work.

\subsection{12.3 The Work of Peter
Scholze}\label{ch2010ngobaochau:the-work-of-peter-scholze}

Peter Scholze (Fields Medal 2018) developed perfectoid spaces and used
them to prove new cases of the Langlands correspondence (the ``local
Langlands correspondence'' for \(\mathrm{GL}_n\)). Ngô's work on the
Fundamental Lemma complements Scholze's work: Ngô deals with the global
aspects of the Langlands program, while Scholze deals with the local
aspects.

\subsection{12.4 The Work of Akshay
Venkatesh}\label{ch2010ngobaochau:the-work-of-akshay-venkatesh}

Akshay Venkatesh (Fields Medal 2018) has applied the Langlands program
to study the distribution of automorphic forms. The tools developed by
Ngô --- especially the stabilized trace formula --- are essential in
this work.

\begin{center}\rule{0.5\linewidth}{0.5pt}\end{center}

\section{13. The Vietnam Institute for Advanced Study in
Mathematics}\label{ch2010ngobaochau:the-vietnam-institute-for-advanced-study-in-mathematics}

\subsection{13.1 Founding}\label{ch2010ngobaochau:founding}

Ngô played a central role in founding the Vietnam Institute for Advanced
Study in Mathematics (VIASM) in 2011. He has served as its Scientific
Director since its inception.

\subsection{13.2 Mission}\label{ch2010ngobaochau:mission}

VIASM's mission is to: - Promote mathematical research in Vietnam -
Train young Vietnamese mathematicians - Host international visitors and
conferences - Connect Vietnamese mathematics to the global mathematical
community

\subsection{13.3 Impact}\label{ch2010ngobaochau:impact}

VIASM has become a major center for mathematical research in Southeast
Asia. Under Ngô's leadership, it has: - Hosted hundreds of international
visitors - Trained thousands of Vietnamese students - Established
research programs in number theory, algebraic geometry, and related
fields

\subsection{13.4 Ngô's Role}\label{ch2010ngobaochau:nguxf4s-role}

Ngô has been closely committed to developing mathematics in Vietnam.
Despite his positions at the University of Chicago and the Institute for
Advanced Study, he has maintained his connection to VIASM and has been
actively involved in its activities.

\begin{center}\rule{0.5\linewidth}{0.5pt}\end{center}

\section{14. Ngô's Broader Mathematical
Influence}\label{ch2010ngobaochau:nguxf4s-broader-mathematical-influence}

\subsection{14.1 Geometric Methods in Representation
Theory}\label{ch2010ngobaochau:geometric-methods-in-representation-theory}

Ngô's work has been instrumental in establishing algebraic geometry as a
central tool in representation theory. Before his work, the Langlands
program was primarily studied using analytic methods (harmonic analysis,
trace formulas). Ngô showed that geometric methods --- especially the
geometry of moduli spaces and perverse sheaves --- could be equally
powerful.

\subsection{14.2 The Hitchin Fibration as a Universal
Tool}\label{ch2010ngobaochau:the-hitchin-fibration-as-a-universal-tool}

The Hitchin fibration, originally introduced in the study of Higgs
bundles and gauge theory, has become a central object in mathematics
thanks in part to Ngô's work. It has been applied to: - The Langlands
program - Mirror symmetry (the Hitchin fibration appears in the
Strominger--Yau--Zaslow conjecture) - Teichmüller theory - Integrable
systems

\subsection{14.3 Influence on Young
Mathematicians}\label{ch2010ngobaochau:influence-on-young-mathematicians}

Ngô has supervised numerous PhD students and postdocs, many of whom have
gone on to become leading mathematicians. His students include: -
\textbf{Bao Le Hung} (University of Chicago): Work on the Langlands
program - \textbf{Duc Ngo} (University of Reading): Work on automorphic
forms

\subsection{14.4 Intergenerational
Transmission}\label{ch2010ngobaochau:intergenerational-transmission}

Ngô's doctoral advisor, Gérard Laumon, was a student of Alexander
Grothendieck at IHÉS. Ngô thus represents the third generation of the
Grothendieck school, and his work carries forward the tradition of using
algebraic geometry to solve problems in number theory and representation
theory.

\begin{center}\rule{0.5\linewidth}{0.5pt}\end{center}

\section{15. Awards and Recognition}\label{ch2010ngobaochau:awards-and-recognition}

\textbf{2010:} - Fields Medal (International Mathematical Union)

\textbf{2007:} - Prix Jacques Herbrand (Académie des Sciences, France)

\textbf{2005:} - Clay Research Award (Clay Mathematics Institute) -
Vietnamese Mathematical Society Award

\textbf{2004:} - Oberwolfach Prize (Oberwolfach Foundation)

\textbf{2003:} - Prix de la Fondation de l'École Polytechnique

\textbf{2002:} - CNRS Bronze Medal

\textbf{1998:} - Prix Michel Monier (Université Paris-Sud)

\textbf{National Honors:} - Received numerous awards from the Vietnamese
government for his contributions to Vietnamese mathematics and
education.

\begin{center}\rule{0.5\linewidth}{0.5pt}\end{center}

\section{16. Selected Major
Publications}\label{ch2010ngobaochau:selected-major-publications}

\begin{enumerate}
\def\labelenumi{\arabic{enumi}.}
\item
 \textbf{``Le lemme fondamental pour les algèbres de Lie.''}
 \emph{Publications Mathématiques de l'IHÉS} 111(1), 1--169, 2010. ---
 \emph{The definitive proof of the Fundamental Lemma for Lie algebras.}
\item
 \textbf{``Le lemme fondamental de Jacquet et Ye en caractéristique
 positive''} (with G. Laumon). \emph{Duke Mathematical Journal} 122(2),
 325--350, 2004. --- \emph{Proof of the Fundamental Lemma for unitary
 groups.}
\item
 \textbf{``Fibration de Hitchin et endoscopie.''} \emph{Inventiones
 Mathematicae} 164(2), 399--453, 2006. --- \emph{Introduces the use of
 the Hitchin fibration for endoscopic theory.}
\item
 \textbf{``Le lemme fondamental pour les groupes unitaires''} (with G.
 Laumon). \emph{Annals of Mathematics} 168(2), 477--573, 2008. ---
 \emph{Comprehensive proof for unitary groups.}
\item
 \textbf{``Decomposition de Springer et le lemme fondamental.''}
 \emph{Journal of the American Mathematical Society} 22(1), 277--325,
 2009. --- \emph{Springer theory and the Fundamental Lemma.}
\item
 \textbf{``Le lemme fondamental: une démonstration géométrique.''}
 \emph{Séminaire Bourbaki}, 2009. --- \emph{Exposition of the geometric
 proof.}
\item
 \textbf{``On the Goormaghtigh equation.''} \emph{Comptes Rendus
 Mathématique} 335(9), 745--748, 2002. --- \emph{Contributions to the
 Goormaghtigh conjecture.}
\item
 \textbf{``On the Sato--Tate conjecture for modular forms''} (with T.
 Barnet-Lamb, D. Geraghty, R. Taylor). --- \emph{Application of
 Langlands program ideas to the Sato--Tate conjecture.}
\end{enumerate}

\begin{center}\rule{0.5\linewidth}{0.5pt}\end{center}

\section{17. Ongoing Work and Future
Directions}\label{ch2010ngobaochau:ongoing-work-and-future-directions}

\subsection{17.1 The Langlands Functoriality
Conjecture}\label{ch2010ngobaochau:the-langlands-functoriality-conjecture}

With the Fundamental Lemma proved, the next major challenge is to fully
establish the Langlands functoriality conjecture --- the transfer of
automorphic representations between arbitrary reductive groups. Ngô's
work provides the foundation for this, but much work remains.

\subsection{17.2 Beyond the Fundamental
Lemma}\label{ch2010ngobaochau:beyond-the-fundamental-lemma}

Ngô's geometric methods --- especially the use of the Hitchin fibration
and the support theorem --- have potential applications beyond the
Fundamental Lemma: - \textbf{The arithmetic fundamental lemma} (a
generalization involving derived categories) - \textbf{The geometric
Langlands program} (where Hitchin fibrations already play a central
role) - \textbf{The theory of Shimura varieties} (used in the study of
modular forms and Galois representations)

\subsection{17.3 Developing Mathematics in
Vietnam}\label{ch2010ngobaochau:developing-mathematics-in-vietnam}

Ngô continues to be closely involved in VIASM and the development of
mathematics in Vietnam. He has been working to: - Increase the number of
Vietnamese mathematicians working internationally - Improve mathematical
education at Vietnamese universities - Establish research programs in
emerging areas of mathematics

\begin{center}\rule{0.5\linewidth}{0.5pt}\end{center}

\section{18. Conclusion}\label{ch2010ngobaochau:conclusion}

Ngô Bào Châu's proof of the Fundamental Lemma represents one of the
greatest achievements in 21st-century mathematics. The Fundamental Lemma
had been the central obstacle to progress in the Langlands program for
over three decades. Ngô's solution --- which combined deep insights from
algebraic geometry (the Hitchin fibration), representation theory
(endoscopy), and number theory (the global-to-local method) --- solved
this problem completely and definitively.

The impact of Ngô's work extends far beyond the Fundamental Lemma
itself: - It completed the stabilization of the trace formula, one of
the central tools in the Langlands program. - It established geometric
methods as essential tools in the study of automorphic forms. - It
opened new directions in the geometric Langlands program and the theory
of Hitchin systems.

The Fields Medal citation --- \emph{``for his proof of the Fundamental
Lemma in the theory of automorphic forms through the introduction of new
algebro-geometric methods''} --- captures the essence of his
contribution. Ngô did not merely prove a difficult conjecture; he
introduced new methods that transformed the field.

Beyond his mathematical achievements, Ngô has been a figure of substantial
importance for Vietnamese mathematics, helping to establish VIASM and
inspiring a new generation of Vietnamese mathematicians. His journey
from a war-torn country to the highest honor in mathematics is a
testament to the power of talent, determination, and the universal
language of mathematics.

\begin{center}\rule{0.5\linewidth}{0.5pt}\end{center}

\begin{center}\rule{0.5\linewidth}{0.5pt}\end{center}

\section*{Exercises}
\addcontentsline{toc}{section}{Exercises}

\begin{exercise} Verify the fundamental lemma for $\\operatorname{GL}_2$ by computing the orbital integrals explicitly. \end{exercise}

\begin{exercise} Show that the trace formula equates a sum over orbital integrals to a sum over character values. \end{exercise}

\begin{exercise} Compute the stable orbital integral for the identity coset in $\\operatorname{GL}_n$. \end{exercise}

\begin{exercise} Prove that the fundamental lemma for $\\operatorname{GL}_n$ follows from the case of $\\operatorname{GL}_2$ using the endoscopic classification. \end{exercise}

\begin{center}\rule{0.5\linewidth}{0.5pt}\end{center}

\section*{Summary}
\addcontentsline{toc}{section}{Summary}

Ng\{o} proved the fundamental lemma, a key conjecture in the Langlands program connecting automorphic forms on different groups. His geometric proof using the Hitchin fibration introduced new techniques from algebraic geometry into number theory, resolving a problem that had been open for over 30 years.

\section{19. References and Further
Reading}\label{ch2010ngobaochau:references-and-further-reading}

\subsection{Official Sources}\label{ch2010ngobaochau:official-sources}

\begin{itemize}
\tightlist
\item
 ICM 2010 Fields Medal Announcement:
 https://www.icm2010.in/prize-winners-2010/fields-medal-ngo-bao-chau
\item
 IMU Fields Medal Citations:
 https://www.mathunion.org/imu-awards/fields-medal/fields-medal-2010
\end{itemize}

\subsection{Biographical and Overview
Articles}\label{ch2010ngobaochau:biographical-and-overview-articles}

\begin{itemize}
\tightlist
\item
 Chau, N. B. ``A Path to the Fundamental Lemma.'' \emph{Notices of the
 AMS} 58(7), 2011.
\item
 University of Chicago News. ``Ngô Bào Châu Receives Fields Medal.''
 August 2010.
\item
 Hartnett, K. ``How a Vietnamese Mathematician Proved the Fundamental
 Lemma.'' \emph{Quanta Magazine}, 2018.
\end{itemize}

\subsection{Key Technical References}\label{ch2010ngobaochau:key-technical-references}

\begin{itemize}
\tightlist
\item
 Ngô, B. C. ``Le lemme fondamental pour les algèbres de Lie.''
 \emph{Publications Mathématiques de l'IHÉS} 111 (2010), 1--169.
\item
 Laumon, G. and Ngô, B. C. ``Le lemme fondamental pour les groupes
 unitaires.'' \emph{Annals of Mathematics} 168 (2008), 477--573.
\item
 Ngô, B. C. ``Fibration de Hitchin et endoscopie.'' \emph{Inventiones
 Mathematicae} 164 (2006), 399--453.
\item
 Langlands, R. ``The Work of Ngô Bào Châu.'' ICM 2010 Laudatio.
\end{itemize}

\subsection{Background Reading}\label{ch2010ngobaochau:background-reading}

\begin{itemize}
\tightlist
\item
 Arthur, J. ``The Fundamental Lemma and the Langlands Program.''
 \emph{Bulletin of the AMS} 51 (2014), 221--249.
\item
 Hitchin, N. ``The Self-Duality Equations on a Riemann Surface.''
 \emph{Proceedings of the London Mathematical Society} 55 (1987),
 59--126.
\item
 Kottwitz, R. ``Stable Trace Formula: Cuspidal Tempered Terms.''
 \emph{Duke Mathematical Journal} 51 (1984), 611--650.
\end{itemize}

\begin{center}\rule{0.5\linewidth}{0.5pt}\end{center}

\section*{See Also}
\addcontentsline{toc}{section}{See Also}

This chapter relates to the following work by other Fields Medalists:
\begin{itemize}
  \item \textbf{\hyperref[ch1966alexandergrothendieck:alexander-grothendieck-fields-medal-1966-a-comprehensive-report]{See Chapter on related work}} --- Ngô's work uses Grothendieck's techniques.
  \item \textbf{\hyperref[ch1978pierrerenedeligne:pierre-renuxe9-deligne-fields-medal-1978-a-comprehensive-report-on-his-work-and-impact]{See Chapter on related work}} --- Deligne's work on the Weil conjectures relates.
\end{itemize}

